\begin{table}[H]
\centering
\captionsetup{justification=centering}
\caption{Teste de diagnóstico do resíduo para a taxa de desocupação - modelos univariado e multivariado}
\label{tab:diagtaxa}
\scalebox{1}{
\renewcommand{\arraystretch}{0.7}
\begin{tabular}{lccc}
\toprule
& \multicolumn{3}{c}{\textbf{Modelo Univariado}} \\
\cmidrule(lr){2-4}
\textbf{Estrato geográfico} & Shapiro-Wilk & Ljung-Box & H \\
\midrule
01 - Belo Horizonte & 0,8601 & 0,9309 & 0,5742 \\
02 - Entorno e Colar Metropolitano de BH & 0,1700 & 0,0319** & 0,6347 \\
03 - Sul de Minas & 0,0512* & 0,0799* & 0,9262 \\
04 - Triângulo Mineiro & 0,6179 & 0,2440 & 0,8540 \\
05 - Zona da Mata & 0,9283 & 0,2884 & 0,6149 \\
06 - Norte de Minas & 0,1454 & 0,0974* & 0,5049 \\
07 - Vale do Rio Doce & 0,8520 & 0,3597 & 0,2295 \\
08 - Central & 0,1325 & 0,1337 & 0,1126 \\
\midrule
& \multicolumn{3}{c}{\textbf{Modelo Multivariado}} \\
\cmidrule(lr){2-4}
\textbf{Estrato geográfico} & Shapiro-Wilk & Ljung-Box & H \\
\midrule
01 - Belo Horizonte & 0,8407 & 0,8420 & 0,5084 \\
02 - Entorno e Colar Metropolitano de BH & 0,6037 & 0,0951* & 0,6204 \\
03 - Sul de Minas & 0,3230 & 0,1650 & 0,9349 \\
04 - Triângulo Mineiro & 0,9237 & 0,2139 & 0,3977 \\
05 - Zona da Mata & 0,0237** & 0,3745 & 0,9795 \\
06 - Norte de Minas & 0,2102 & 0,4442 & 0,3363 \\
07 - Vale do Rio Doce & 0,4342 & 0,4100 & 0,2346 \\
08 - Central & 0,6965 & 0,6044 & 0,4538 \\
\bottomrule
\end{tabular}}
\fonte{Elaboração própria, com base nos dados da PNAD Contínua. Nota: * representa rejeição de \(H_0\) a 10\%, ** a 5\% e *** a 1\%. As oito primeiras observações foram descartadas, em razão do período de estabilização do filtro de Kalman.}
\end{table}
